\documentclass[10pt,letterpaper]{article}

\usepackage{compose_arxiv}

\newcommand{\method}{\textsc{Compose}\xspace}
\newcommand{\rgm}{\textsc{RGM}\xspace}
\newcommand{\X}{\mathcal{X}}
\newcommand{\A}{\mathcal{A}}
\newcommand{\G}{\mathcal{G}}
\newcommand{\T}{\mathsf{T}}
\newcommand{\E}{\mathbb{E}}

\hypersetup{
  pdftitle={COMPOSE: Generator Matching over Stochastic Rewrites for Trans-Dimensional Molecular Generation and Control},
  pdfauthor={Anonymous Authors},
  pdfsubject={Executable stochastic graph rewrites for molecular generation and control},
  pdfkeywords={molecular generation, graph rewriting, generator matching, trans-dimensional processes, stochastic control},
  pdfcreator={LaTeX with hyperref}
}

\begin{document}

\begin{composefrontmatter}
\composetitle{COMPOSE: Generator Matching over Stochastic Rewrites for
Trans-Dimensional Molecular Generation and Control}
\composeauthors{Anonymous Authors}
\composeabstract{%
Generative models for molecular graphs commonly learn to reverse a convenient
corruption on an ambient representation. Such transitions need not be
executable chemical edits, and intermediate states need not be complete
molecules. We introduce \emph{Rewrite Generator Matching} (RGM), a framework
that instead learns a time-inhomogeneous marked jump process over the
state-dependent fiber of legal executable rewrites. A total hazard determines
when an event occurs and a conditional mark law determines which rewrite
fires. We instantiate RGM as COMPOSE on disjoint cardinality strata, including
a distinguished null source that is not a molecule. Its marks create, delete,
restate, reconnect, and cyclize molecular structure, while the executor keeps
every non-null committed state inside a declared space of connected,
valence-valid molecules. Because different marks can execute to the same
molecule, we define the molecular transition law as the pushforward of the
learned mark law through execution and canonical molecular identity. This
successor kernel is invariant to probability-preserving refinements within a
successor fiber and is the object on which finite-horizon conditioning, hard
support constraints, and dynamic retargeting act. Exact control is claimed
only on bounded enumerable state spaces; learned values approximate it at
molecular scale. Source-conditioned editing and timed de novo generation
require separate checkpoints and inference regimes.}
\end{composefrontmatter}

\section{Introduction}
\label{sec:introduction}

Much of modern generative modeling begins by choosing a tractable corruption
process and learning to reverse it. Gaussian noise, categorical replacement,
and continuous-time discrete corruption have produced powerful general
formalisms \citep{ho2020ddpm,austin2021d3pm,campbell2022ctdd}; their graph
instantiations place probability on represented node and edge states
\citep{vignac2023digress,qin2025defog}. For molecular design, however, the
transition mechanism is more than mathematical scaffolding. A useful
intermediate should be a molecule on which a property model, constraint, or
chemist can act, and a transition should specify an operation that can actually
be executed. Generic corruptions do not promise either property: intermediate
graphs need not belong to the molecular state space, and a tensor change need
not correspond to a chemically meaningful edit.

We ask whether the generative process itself can instead be a stochastic
program of executable rewrites. The premise is simple. When a domain supplies
typed rewrite rules, state-dependent application conditions, and a
deterministic executor, those objects can define the infinitesimal events of a
learned jump process. The model then learns transport through operations that
already possess domain meaning, rather than learning to undo a process chosen
only for analytical convenience. This changes the role of an edit vocabulary:
it is not merely a proposal mechanism attached to an optimizer, but the legal
support of the generative process.

\paragraph{Generator Matching over executable marks.}
We formalize this idea as \emph{Rewrite Generator Matching} (\rgm). At state
$x$, the rewrite system exposes a finite, state-dependent action fiber
$\A(x)$. A mark $a\in\A(x)$ is a complete executable operation---a rule,
concrete match or address, and typed payload---and the deterministic executor
$\T(x,a)$ returns its successor. The model factorizes the marked rate as
\begin{equation}
  \lambda_\theta(a\mid x,t)
  =
  \Lambda_\theta(x,t)\,p_\theta(a\mid x,t),
  \qquad
  a\in\A(x),
\end{equation}
where $\Lambda_\theta$ is the total event hazard and $p_\theta$ is the
conditional mark law. Its generator acts on test functions by
$\sum_{a\in\A(x)}\lambda_\theta(a\mid x,t)
[f(\T(x,a))-f(x)]$. Generator Matching supplies the rate-fitting principle
\citep{holderrieth2025generator}; the executable rewrite system supplies the
state space, event fibers, legal support, and endpoint-conditioned programs
from which teacher events are drawn. The learned object is therefore a marked
stochastic rewrite process, not an unconstrained prediction followed by
repair.

\paragraph{A genuinely trans-dimensional molecular process.}
\method instantiates \rgm on the disjoint union
\[
  \X=\bigsqcup_{n=0}^{N}\X_n,
  \qquad
  \X_0=\{\varnothing\}.
\]
For $n\geq 1$, $\X_n$ contains supported, connected, valence-valid molecular
graphs with $n$ active atoms. The distinguished null state $\varnothing$ is
\emph{not} a molecule; it is a reversible source connected to supported
one-atom states. Production tensors use persistent padded slots, but padding is
only a coordinate system. Atom birth and death change the active graph and
therefore move between different cardinality strata.

The molecular marks create and delete atoms, restate atom type and valence,
change bond order, reroute attached subgraphs, and close or open cycles. The
legal fiber depends on the current
molecule, and every committed action is executed before the successor is
accepted. Consequently, every \emph{non-null} committed state is a complete
molecule and every transition is an executable rewrite. This is the structural
reason a trajectory can support mid-course property evaluation, protected
substructures, interruption, branching, size correction, and topology
redesign. It does not imply synthetic accessibility, and the declared operator
set need not be one-step inverse-closed for every operation.

\paragraph{The molecular law is a pushforward.}
Different executable marks can yield the same molecular successor because of
symmetry, persistent-slot aliases, operand ordering, or alternative internal
refinements. For a canonical molecular successor $y$, define its successor
fiber
\[
  \G_y(x)=\{a\in\A(x):\T(x,a)\cong y\}.
\]
The state-level molecular law is not any one representative mark. It is the
pushforward
\[
  P_\theta(y\mid x,t)
  =
  \sum_{a\in\G_y(x)}p_\theta(a\mid x,t),
\]
with the embedded-chain convention treating molecular self or virtual events
consistently. This distinction is essential. Refining or coarsening marks
within one successor fiber, while preserving their aggregate mass, leaves the
molecular process unchanged. It also leaves any controller defined on
canonical successors unchanged. A mark-level power tilt, top-$k$ rule, nucleus
filter, or family-wise bias generally lacks that invariance. We therefore
control and evaluate the process at the molecular-successor level even when
the internal model is factorized over marks.

This transition-level aggregation is closely related to pushforward Generator
Matching for deterministic images of latent processes
\citep{billera2026latent}. That theory is the closest generic precedent; here
aliases are executable event marks, the map
$a\mapsto\operatorname{canon}(\T(x,a))$ depends on the current molecule, and
the aggregated kernel is retained as the molecular control interface.

\paragraph{Control follows from the kernel.}
Because every productive edge in $P_\theta$ is legal and every non-null state
admits molecular evaluation, the induced kernel is an intervention substrate.
Hard path constraints restrict its support; soft objectives reweight supported
successors; and a remaining-budget value permits continuation after an
objective changes. For a fixed kernel $P$, terminal desirability $g$, and
remaining budget $b$, the finite-horizon value
$h_b(x)=\E_P[g(X_b)\mid X_0=x]$ gives the Doob-controlled law
$P_b^g(y\mid x)=P(y\mid x)h_{b-1}(y)/h_b(x)$ whenever $h_b(x)>0$.
The transform never creates an illegal edge, although hard conditioning may
prune base support. We claim exactness only on bounded state spaces that can be
enumerated and audited. At molecular scale, a learned value function is an
approximation and must be calibrated as such.

\paragraph{Two checkpoints, two inference regimes.}
The framework supports two related but distinct uses. Source-conditioned
editing starts from a supplied molecule and follows a fixed-budget embedded
molecular jump chain; it does not use the magnitude of the continuous-time
hazard. Unconditional de novo generation starts from its own source
distribution, includes the reversible null source, and uses a timed
Gillespie-style CTMC \citep{gillespie1977}. These regimes require separate data
mixtures, hazard contracts, checkpoints, and evaluations. An editing
checkpoint is not evidence for unconditional generation, even though both
instantiate the same rewrite framework.

\paragraph{Training and claim boundary.}
The trained random variable and the induced molecular random variable must not
be conflated. The completed 16,000-step editing run optimized the implemented
mark-level Generator-Matching objective, with successor aggregation only where
that implementation explicitly provided it, and remains a diagnostic rather
than a selected paper model. A repaired protocol that optimizes
canonical-successor likelihood is a distinct training protocol; it does not
retroactively make that run successor-trained. In either case, the scientific
definition of the molecular kernel is execution followed by canonical
pushforward, not the particular segmented algorithm used to compute it.
The conditional-loss result of \citet{billera2026latent} applies when the model
and loss parameterize the image-space generator under its assumptions; it does
not turn a selected-mark objective into a successor objective.

\paragraph{Contributions.}
One causal chain organizes the paper: executable state-dependent rewrites lead
to Generator Matching, which yields a trans-dimensional molecular process;
canonical pushforward turns that marked process into a representation-independent
molecular kernel; and that kernel supports exact finite-state and approximate
large-scale control. Our contributions are:

\begin{enumerate}
  \item \textbf{Generator Matching over stochastic rewrites.} We formulate a
  learned marked jump process over state-dependent executable action fibers,
  separating total event hazard from the conditional identity and location of
  a rewrite.

  \item \textbf{COMPOSE and the canonical molecular-successor law.} We
  instantiate the framework with molecular birth, death, relabeling,
  reconnection, and compositional cycle operations across cardinality strata,
  and operationalize its molecular control kernel by aggregating executable
  marks through canonical molecular identity.

  \item \textbf{Generation and control on one legal molecular substrate.} We
  provide a common successor-kernel interface for exact finite-horizon
  conditioning on bounded enumerable spaces and for learned approximations to
  dynamic retargeting, Pareto branching, and pathwise constraints at molecular
  scale, while keeping fixed-budget editing and timed de novo generation
  explicitly separate.
\end{enumerate}

\paragraph{Relation to prior work.}
None of atom insertion or deletion, valid molecular editing, graph grammars,
Generator Matching, variable-dimension generation, Pareto optimization, or
Doob control is individually new. Discrete diffusion and flow models provide
general graph-generation mechanisms
\citep{austin2021d3pm,campbell2022ctdd,vignac2023digress,qin2025defog};
recent edit-based processes make insertion and deletion first-class
\citep{havasi2025editflows,ninniri2025griddd}. Molecular grammars define
structured legal support \citep{kajino2019mhg,guo2022grammar}, while stochastic
graph rewriting supplies rule applications and CTMC semantics
\citep{danos2015thermodynamic,behr2021rewriting}. Molecular editors and
optimization systems already traverse valid molecular neighborhoods
\citep{xie2021mars,fu2021mimosa,gao2022pmo}. The contribution here is their
particular formal composition: a Generator-Matched executable rewrite process
with trans-dimensional molecular semantics, a canonical successor
pushforward, and successor-level exact and dynamic control.
Pushforward Generator Matching supplies the closest general result for
deterministic images of latent processes \citep{billera2026latent}; we do not
claim that generic theorem as a contribution.

\bibliographystyle{plainnat}
\bibliography{references}

\end{document}
